\title{LongRoPE: Extending LLM Context Window Beyond 2 Million Tokens}

\begin{abstract}
Expanding the context window is highly sought after in large language models (LLMs). Nevertheless, contemporary context window extensions typically max out at roughly 128k tokens, owing to significant fine-tuning expenses, limited availability of lengthy textual data, and disruptive effects caused by unfamiliar token positions.

This study presents LongRoPE, a method that, for the first time, expands the context window of pre-trained LLMs to an extraordinary \textbf{2048k} tokens, requiring as few as 1k fine-tuning iterations on sequences of up to 256k tokens, while preserving performance at the original, shorter context length. This capability arises from three main advancements: (i) we discover and leverage two distinct forms of non-uniformities in positional interpolation via a computationally efficient search procedure, leading to superior initialization for fine-tuning and facilitating up to an 8$\times$ increase in context length even without additional fine-tuning; (ii) we propose a staged extension approach in which the model is first fine-tuned on sequences of 256k tokens, followed by a second application of positional interpolation on the fine-tuned model, resulting in a context window of 2048k tokens; (iii) we further adjust LongRoPE at the 8k token length to fully regain performance on short sequences. Comprehensive evaluations using LLaMA2 and Mistral, covering multiple benchmark tasks, confirm the robustness and efficacy of our technique. Models modified through LongRoPE largely retain their original architecture, introducing only minimal changes to the positional embeddings, and remain compatible with most established optimizations. Source code will be released at \texttt{https://github.com/microsoft/LongRoPE}

\section{Introduction}

Although Large Language Models (LLMs) have achieved impressive results across diverse tasks~\cite{gpt4,llama2}, their ability is constrained by a finite context window, such as the 4096-token maximum in LLaMA2~\cite{llama2}. When inputs exceed this limit, the models exhibit degraded performance, as they encounter positional configurations outside their training distribution. These limitations are particularly problematic in scenarios requiring in-context learning with substantial numbers of examples~\cite{Huang2023FewerIM}, as well as in the deployment of LLM-based agents~\cite{park2023generative,madaan2023selfrefine}.

Recent studies demonstrate that the context window of a pre-trained LLM can be expanded to approximately 128k tokens through fine-tuning with longer sequences~\cite{longlora,pi,yarn,zhang2024soaring,liu2023scaling}. However, advancing beyond this limit presents three significant challenges. Firstly, the incorporation of previously unseen position indices introduces many anomalous values, which disrupt the model’s training distribution and hinder the fine-tuning process~\cite{pi}. This problem is exacerbated in cases where the context window is enlarged from 4k to over 1000k, thereby generating more than 90\% new positional tokens. Secondly, successful fine-tuning is contingent upon the availability of suitably lengthy texts, yet existing datasets rarely contain texts surpassing 1000k tokens. Additionally, processing these extremely long sequences demands substantial computational effort, resulting in impractically high training durations and considerable GPU usage. Thirdly, as context windows are scaled to extreme lengths, the attention mechanism is diluted—tokens must attend to a vast number of positions—which undermines performance when the model is tested on shorter contexts~\cite{pi}.

To address the initial challenge, a commonly used strategy involves interpolating RoPE positional embeddings~\cite{rope,pi} by scaling novel position indices to fit within the original pretrained limits, as illustrated in Fig.\ref{fig:overview}. The Position Interpolation (PI) method~\cite{pi} adjusts RoPE’s rotary angles through linear interpolation based on the window extension ratio. NTK~\cite{ntk,dynamicntk} proposes a non-uniform approach to interpolation and extrapolation among the different RoPE dimensions. YaRN~\cite{yarn} groups RoPE dimensions into three categories according to their frequencies, employing extrapolation, NTK-based, and linear interpolation techniques for each respective group. Despite these methods, positional embeddings in the Transformer model are characterized by \emph{complex, non-uniform information entropy}. Current methods fail to capitalize on these subtle, non-uniform features, which results in the loss of information and restricts the ability to extend the context window.

Section~\ref{sec:analysis} presents empirical evidence for two major observations: \textbf{(1)} Successful positional interpolation must take into account two distinct types of non-uniformity—namely, the heterogeneity across RoPE dimensions as well as across different token positions. Both reduced RoPE dimensions and earlier token positions warrant limited interpolation, though the ideal approach is contingent upon the specific target length to be extrapolated. \textbf{(2)} Incorporating these non-uniform aspects into positional interpolation ensures more effective preservation of the original RoPE's information, particularly in regard to crucial dimensions and early token positions. This approach helps mitigate the information loss typically induced by positional interpolation, which in turn supplies superior initial weights for subsequent fine-tuning. Additionally, this technique supports up to 8$\times$ extension even without additional fine-tuning.

Building upon these insights, we present \textbf{LongRoPE}, a robust approach for extending the context window of LLMs to lengths exceeding 2 \emph{million} tokens. The design of LongRoPE relies on three fundamental advancements. Firstly, {LongRoPE} leverages the complex non-uniformities across multiple dimensions in positional interpolation. It determines optimal rescale factors for RoPE’s rotation angles in each dimension, tailored to token position. Since the search space for suitable rescale factors grows exponentially with higher extension ratios, {LongRoPE} employs an evolutionary search algorithm that incorporates two optimization strategies to improve search performance. Fig.~\ref{fig:overview} illustrates the outcome of these optimized, rescaled RoPE parameters.

Subsequently, {LongRoPE} employs an efficient and gradual extension approach to attain a 2048k context window, circumventing the necessity of fine-tuning on exceptionally lengthy texts, which are rare and difficult to procure. The procedure initiates by identifying a 256k sequence length on the pre-trained LLM, followed by fine-tuning at this specific length. With our method of non-uniform positional interpolation—enabling an 8$\times$ increase in non-fine-tuning circumstances—a secondary search for optimal RoPE rescale parameters is performed on the extended, fine-tuned LLM. Through this process, the 2048k context window for LLaMA2 and Mistral~\cite{mistral} is ultimately realized.

In order to prevent a decline in performance within the initial, smaller context window, {LongRoPE} persistently refines the RoPE rescale parameters on the enlarged LLM. Following the same methodology applied when increasing the context window from 256k to 2048k, we perform a reduction to 4k and 8k windows on the LLM that has been fine-tuned at 256k, utilizing our search algorithm to reduce positional interpolation. At inference time, when the input sequence is shorter than 8k, RoPE is updated using the optimized rescale factors identified by the search process.

A comprehensive set of experiments involving multiple LLMs and a range of long-context tasks substantiates the efficacy of our proposed approach. The results indicate that LongRoPE consistently sustains low perplexity when evaluated at lengths between 4k and 2048k, attains passkey retrieval accuracy exceeding 90\%, and achieves performance on par with baseline models on established benchmarks that operate within a 4096-token context window. LongRoPE is compatible with all LLMs that utilize RoPE embeddings. Our codebase, along with the LongRoPE-2048k models, will be made publicly available.

\section{Non-uniformity in Positional Interpolation}
\label{sec:analysis}

\subsection{Preliminary}

Transformer models require explicit positional information, often in the form of position embedding, to represent the order of input tokens. Our work focuses on the RoPE~\cite{rope} position embedding, which is widely used in recent  LLMs. For a token at position index $n$, its corresponding RoPE encoding can be simplified as follows:

\begin{equation}
		\fontsize{7.8}{7.8} \selectfont
	\label{eq:rope}
	\left[ cos(n\theta_0), sin(n\theta_0), cos(n\theta_1), \cdot\cdot\cdot, cos(n\theta_{d/2-1}), sin(n\theta_{d/2-1}) \right]   
\end{equation}

Here, $d$ refers to the embedding dimension, $n\theta_i$ denotes the rotary angle for the token located at position $n$, and $\theta_i=\theta^{-2i/d}$ specifies the corresponding rotational frequency. RoPE typically employs a default value of 10000 for the base $\theta$.

\textbf{\emph{Context window extension ratio $s$} and positional interpolation}. The extension ratio $s$ is defined as the proportion between the extended context length $L'$ and the initial context length $L$, expressed as $s=\frac{L'}{L}$.

To extend the context window from $L$ to $L'$, current positional interpolation methods suggest downscaling rotation frequencies $\theta_i$ based on extension ratio $s$. Let $\beta=\theta^{2/d}$, and $\lambda$ denote the actual rescale factor related to $s$, we   unify these positional interpolation methods as follows:

\begin{equation}	
	\fontsize{7.5}{7.5} \selectfont
	\label{eq:unified_rope}
	\begin{aligned}
		\left[cos\left(\frac{n}{\lambda(\beta)^0}\right), sin \left(\frac{n}{\lambda(\beta)^0}\right), cos\left(\frac{n}{\lambda(\beta)^1}\right), \cdot\cdot\cdot,  sin \left(\frac{n}{\lambda(\beta)^{d/2-1}}\right) \right]   
	\end{aligned}
\end{equation}

\textbf{Linear positional interpolation (PI)}. In PI~\cite{pi}, positions are interpolated linearly, confined to the range of indices used during pre-training. To achieve a sequence extension by a factor $s$, every rotation angle is uniformly scaled down by $\lambda = s$ over each RoPE dimension. This approach compresses positional information, resulting in densely packed representations that impede the model’s discrimination between nearby tokens. Consequently, PI demonstrates notable performance drops when applied to significant sequence extensions.

\textbf{NTK-based interpolation and extrapolation}. ~\cite{ntk,dynamicntk} reinterpret RoPE through the framework of information encoding and utilize insights from Neural Tangent Kernel (NTK) theory~\cite{jacot2018neural,tancik2020fourier}. To address the problem of crowded positions inherent in PI, they propose to balance the interpolation effects over the different RoPE dimensions. This entails applying smaller scaling factors to the lower (i.e., higher frequency) dimensions, while giving larger scaling to the higher (i.e., lower frequency) dimensions. This approach facilitates both positional interpolation and extrapolation, parametrized as $\lambda=s^{i}$. In dynamic NTK~\cite{dynamicntk}, the extension ratio is modulated dynamically depending on the input sequence length. In contrast to PI, which depends on fine-tuning to extend context windows, NTK-based approaches can enlarge the context window without requiring fine-tuning, though their maximum practical extension is generally capped at about 4$\times$.

\textbf{YaRN}~\cite{yarn} presents a substantial advancement in the effectiveness of positional interpolation. It achieves this by organizing RoPE dimensions into three distinct groups, categorized according to their frequency characteristics, with each group subjected to a separate interpolation technique. Specifically, the highest frequency dimensions are extrapolated ($\lambda$=1), while the lowest frequencies adopt linear interpolation (PI). Dimensions at intermediate frequencies are treated using the NTK method. Central to YaRN’s design is the dimension grouping process; however, this relies primarily on empirical, manual selection, which has the potential to yield less than optimal outcomes when applied to novel LLM architectures.

\subsection{Study on Non-uniform Positional Interpolation}

Drawing upon the approaches utilized in NTK and YaRN, we observe that their improvements stem largely from leveraging non-linearity, with particular attention given to varying frequencies among RoPE dimensions to achieve more tailored interpolation and extrapolation. Despite these advancements, prevailing non-linear functions predominantly depend on manual, human-crafted heuristics. This observation prompts us to consider two key issues: (1) Does the existing method for positional interpolation represent the best possible solution? (2) Might there be yet unidentified forms of non-linearity worth exploring?

In order to address these questions, we employ evolution search (refer to Sec.~\ref{sec:method}) to identify improved non-uniform positional interpolations tailored for LLaMA2-7B. The search process relies on perplexity as a metric, evaluating 5 randomly selected samples from the PG19~\cite{pg19} validation dataset. Our empirical investigation yields several important insights.

\textbf{Finding 1}: \textit{RoPE dimensions exhibit substantial non-uniformities, which are not effectively handled by current positional interpolation methods.}

We determine the optimal $\lambda$ value for every RoPE dimension as described in Eq.~\ref{eq:unified_rope}. Table~\ref{tbl:exp1} presents the perplexity outcomes of LLaMA2-7B across various approaches when evaluated on PG19 and Proof-pile~\cite{proof-pile} test sets, all without any model fine-tuning. The parameters identified through our search yield considerable perplexity reductions, indicating that both linear (PI) and non-uniform (Dynamic-NTK and YaRN) interpolation strategies currently in use are not ideal. Interestingly, YaRN exhibits inferior performance compared to PI and NTK on PG19, failing to attain the designated context window size in LLMs that have not been fine-tuned. As an illustration, with an 8k context size, YaRN’s perplexity surges noticeably after reaching 7k tokens.

Our investigation revealed that the rescaled coefficients $\lambda$ from Eq.~\ref{eq:unified_rope} are not homogeneous; this distinguishes them from the constant scaling $s$ applied in PI, NTK's analytic approach, or YaRN's group-wise scaling strategy. This introduction of heterogeneity in the scaling factors yields notable gains in LLaMA2's language modeling performance—measured by perplexity—for context window sizes of 8k and 16k, and this improvement is achieved without any model fine-tuning. The main reason for this benefit is that the modified positional embedding remains closely aligned with the original RoPE structure, notably for key embedding dimensions. Consequently, this helps large language models better differentiate between proximate token positions, alleviating confusion at greater context lengths.

\textbf{Finding 2}: \textit{RoPE for the initial tokens in the input sequence should be extrapolated with less interpolation.} 

We propose that for the first $\hat{n}$ tokens of an input sequence, position interpolation through RoPE should be minimized. The rationale is that these initial tokens typically attract high attention scores and thus play a pivotal role in attention processing, a phenomenon previously noted in Streaming LLM~\cite{streamingllm} and LM-Infinite~\cite{han2023lminfinite}. To examine this, we expand the context window to 8k and 16k by applying PI and NTK, varying the number of leading tokens $\hat n$ (ranging from 0, 2, ..., 256) that are exempt from interpolation. Setting $\hat n$=0 defaults back to the standard PI and NTK methods. Table~\ref{tbl:tokenposition} presents two significant findings: \textbf{(1)} exempting the earliest tokens from position interpolation leads to measurable gains in performance, and \textbf{(2)} the optimal count for $\hat n$ is contingent upon the chosen length of context window extension.

\textbf{Finding 3}: \textit{Non-uniform positional interpolation effectively extends LLM context window in both fine-tuning and non-fine-tuning settings.} 

Although our application of searched non-uniform position interpolation markedly enhances extension capabilities at 8k and 16k context lengths without the need for fine-tuning, handling even longer sequences necessitates a fine-tuning phase. Therefore, we perform fine-tuning on LLaMA2-7B utilizing our optimized RoPE method to support a 64k context window (further details are provided in the Appendix). As indicated in Table~\ref{tbl:finetune}, our approach consistently surpasses both PI and YaRN, regardless of whether LLaMA2-7B is fine-tuned. This superior performance is attributable to our strategy of leveraging non-uniform positional interpolation, which reduces information degradation and offers more favorable starting conditions for subsequent fine-tuning.

\textbf{Summary}. Our research identifies two sources of non-uniformity: differences across RoPE dimensions and variation in token positions. Effectively incorporating these non-uniformities in positional interpolation leads to substantial gains in extending the context length of LLMs.

\section{LongRoPE}

\label{sec:method}
Building upon these insights, we propose LongRoPE, which incorporates a streamlined search algorithm to leverage the two identified non-uniformities. This approach subsequently enables the expansion of the LLM context window to more than 2 million tokens.

\subsection{Problem Formulation}

These two irregularities expand the solution space significantly and create intricate challenges in the optimization process. To mitigate this, we reinterpret the multidimensional non-uniform position interpolation optimization task as a search problem.

For a LLM targeting a  context window size of $L'$ and lengthy input documents $\mathbf{X}$, where each $\mathbf{x}\in\mathbf{X}$ surpasses $L'$ in token length, we denote the original rotary angle of the $i^{th}$ dimension in RoPE embedding at token position $n$ as  $\frac{n}{\beta^i}$. The optimization problem is then formulated as follows:

\begin{equation}
\begin{gathered}
		\label{eq:problem}

\underset {\mathbf{x}\in\mathbf{X}; \, |\mathbf{x}|\ge L'}{ \operatorname {arg\,min} } \,  \mathcal{L}\, \left( \text{LLM} (\text{RoPE}, \mathbf{X})\right), \;\text{where}
\\
\scalebox{0.93}{
$\underset {\begin{subarray}{l} i=0,\dots,\frac{d}{2}-1;\\ n\in[ 0,|\mathbf{x}|);\end{subarray}}{\text{RoPE}_(n)}= {\left[\dots, \cos\left(\mathbb{I}(\hat{\lambda}_{i}, \hat{n}) \frac{n}{\beta^i}\right), \sin\left(\mathbb{I}(\hat{\lambda}_{i}, \hat{n}) \frac{n}{\beta^i}\right),\dots \right]}$}
\\
	\text{where} \quad \mathbb{I}(\hat{\lambda}_{i},\hat{n})=\begin{cases}
	1 &\;\mbox{$n<\hat{n}$}\\
{\frac{1}{\lambda}_{i}} &\;\mbox{$n\geq\hat{n}$}
\end{cases}

In this context, a collection of rescale factors, $\mathbb{I}(\hat{\lambda}_{i},\hat{n})$, is introduced to address two distinct types of non-uniformity. Here, $\hat{\lambda}_i$ corresponds to dimension-level non-uniformity within RoPE, while $\hat{n}$ indicates irregularities related to token position. More precisely, $\mathbb{I}(\hat{\lambda}_{i},\hat{n})$ serves to adjust the rotation angle for the $i^{th}$ RoPE dimension, where the parameter $\hat\lambda_{i}$ provides the adjustment and $\hat{n}$ signifies the threshold for token position. For the first $\hat{n}-1$ positions, the scaling factor $\hat\lambda_{i}$ remains inactive, employing the standard RoPE angle $\frac{n}{\beta^i}$. Once tokens reach positions $n\ge\hat{n}$, the rescale factor becomes operative.

With a predetermined context window length of $L'$, the task involves determining the optimal rescale coefficients ($\mathbb{I}(\hat{\lambda}_{0},\hat{n})$, $\mathbb{I}(\hat{\lambda}_{1},\hat{n})$, ...,$\mathbb{I}(\hat{\lambda}_{i},\hat{n})$...) across the $1^{st}$ through $d^{th}$ dimensions of the RoPE mechanism. Consequently, employing these rescaled $\text{RoPE}$ parameters within the target $\text{LLM}$ should yield the lowest possible next-token prediction loss, $\mathcal{L}$ (corresponding to perplexity), on sequences $\mathbf{X}$ containing $L'$ tokens.

\subsection{Searching the Non-uniform Position Interpolation}

To address the challenge presented in Eq.~\ref{eq:problem}, we propose a straightforward and efficient approach that identifies the best RoPE rescale factors, allowing for comprehensive utilization of the multidimensional non-uniform characteristics within position embedding.

\textbf{Search space}. We construct an extensive search space in order to account for the two types of non-uniformity. Table~\ref{tbl:searchspace} presents the full structure of this space. In particular, we permit tuning of an individual rescale factor for every dimension in the RoPE embedding. To facilitate straightforward search space formulation, the search is performed over $\lambda_i$ and $\hat{n}$, rather than directly on $\mathbb{I}(\hat{\lambda}_i, \hat{n})$, with $\hat{\lambda}_{i}$ defined as $1/\lambda_i$. According to Table~\ref{tbl:searchspace}, the permissible range for $\lambda_i$ begins at 1.0 (indicating direct extrapolation) and extends up to $s \times 1.25$ (representing broader interpolation compared to PI) in increments of 0.01, where $s$ is the ratio for extending the target context window.

The parameter $\hat{n}$ determines how many of the starting token positions preserve the original RoPE embedding without employing position interpolation. In practical experiments, $\hat{n}$ is selected from the set \{0, 1, 2, 4, 8, 12, 16, 20, 24, 28, 32, 64, 128, 256\}. For the case where $\hat{n}=0$, rescale factors identified through search are applied across every token position.

\textbf{Evolution-based search}. The search space outlined in Table~\ref{tbl:searchspace} encompasses a vast array of potential positional interpolation configurations, which makes comprehensive exploration highly demanding. For instance, when considering an extension ratio of $s=4\times$, the number of possibilities becomes $400^{128/2}\times14 = 4\times10^{167}$. As this ratio increases, the size of the search space grows at an exponential rate. To efficiently navigate this immense space, we employ evolution search~\cite{guo2020single}, supplemented by two optimization strategies designed to significantly enhance the speed and effectiveness of the search. The full methodology is detailed in Algorithm~\ref{alg:search}.

\textit{Optimized initial population generation}. Rather than creating an initial population of $P$ rescale factors through entirely random selection, we explicitly include the three RoPE rescale factors associated with PI, NTK, and YaRN in the initial group. The other $P$-3 members of the population are formed by randomly mutating the values of these three rescale factors, each with probability $p$.

\textit{Monotonically non-decreasing constraint}. Following the creation of the initial population, LLM perplexity is evaluated for each candidate. The designated RoPE rescale factors are assigned to the target LLM, and perplexity is measured on the input $\mathbf{X}$. The best k candidates are selected as parents for the evolutionary process. Nonetheless, due to the extensive search space, basic mutation and crossover strategies may navigate toward suboptimal solutions, resulting in redundant perplexity assessments. This inefficiency is amplified in scenarios with large $L'$, where the computation of each perplexity is particularly resource-intensive.

To mitigate this issue, we enforce a monotonicity constraint whereby the sampled RoPE rescaling factors must satisfy $\lambda_i\le \lambda_{i+1}$. Only those RoPE configurations adhering to this restriction are considered during the LLM perplexity assessment, thereby greatly streamlining the search process. Specifically, the rescaling factor $\lambda_i$ is required to increase monotonically across RoPE dimensions ($i$ = 0,...,63). This monotonicity across dimensions is grounded in NTK theory~\cite{jacot2018neural,tancik2020fourier,ntk}, which posits that lower-indexed dimensions, corresponding to higher frequencies, demand less interpolation (smaller $\lambda_i$), while dimensions associated with lower frequencies can accommodate more interpolation (larger $\lambda_i$).

\begin{algorithm}[t]
	\small
	\caption{The search algorithm}
	\textbf{Input:} target LLM, input samples $\mathbf{X}$, population size $P$, mutation size $N_1$, crossover size $N_2$, max iterations $\mathcal{T}$, mutate probability $p$\\

	\begin{algorithmic}[1]
		\label{alg:search}
		\STATE Top-k $\leftarrow$ $\phi$;	
		\STATE $\text{P}_0 \leftarrow$ \textit{Initialize\_population\_with\_optimization} ($P$, $\mathbf{X}$, $p$); 
		\FOR{$i = 1$ to $\mathcal{T}$}
		\STATE \textit{Compute\_perplexity} (LLM, $\text{P}_{i-1}$, $\mathbf{X}$);
		\STATE Top-k $\leftarrow$ \textit{Update\_Topk} (Top-k, $\text{P}_{i-1}$);
		\STATE $\text{P}_{mutation}$ $\leftarrow$ \textit{Mutation\_with\_mono\_constraint} (Top-k, $N_1$, $p$); 
		\STATE $\text{P}_{crossover}$ $\leftarrow$ \textit{Crossover\_with\_mono\_constraint} (Top-k, $N_2$);
		\STATE $\text{P}_i$ = $\text{P}_{mutation}$ $\cup$ $\text{P}_{crossover}$ $\cup$ Top-k;
		\ENDFOR
		\STATE Return the member in Top-k exhibiting the lowest perplexity;
	\end{algorithmic}
\end{algorithm}

\textbf{8$\times$ extension without fine-tuning}. By employing an evolutionary search strategy, our approach successfully discovers optimal, non-uniform RoPE rescale coefficients, maintaining essential dimension and position information to reduce information degradation during interpolation. As illustrated in Fig.\ref{fig:non-ft}, this enables us to expand the context length for LLaMA2 from 4k to 32k without any fine-tuning. In comparison, current techniques like PI, as well as non-uniform NTK and YaRN, experience a rapid increase in perplexity when the context window is extended by a factor of 2$\times$.

\subsection{Extending LLM Context Window to 2048K}
\label{sec:extendto2048k}

\textbf{Progressive extension to 2048k}. In this section, we present our strategy for expanding the context window of pre-trained LLMs beyond the standard 4k limit, reaching up to over 2048k. Our approach utilizes non-uniform positional interpolation, which allows for an 8$\times$ context increase without the need for additional fine-tuning. However, for more substantial expansions, such as scaling by 512$\times$, fine-tuning becomes unavoidable. One possible approach involves identifying suitable RoPE rescaled factors corresponding to the 2048k target and subsequently performing fine-tuning. Nonetheless, this process poses significant obstacles, as it demands vast computational resources and, according to our observations, fine-tuning LLMs at such high extension ratios remains highly challenging (refer to Appendix for further discussion).

Fortunately, LongRoPE proves to be successful when applied to both the base and fine-tuned variants of the extended LLM. To this end, we propose a progressive and resource-efficient approach that reaches the desired 2048k target through only 1k fine-tuning iterations, utilizing a training sequence length of up to 256k.

$\Diamond$ \textit{LongRoPE search for extending pre-trained LLM to 256k.} Using LLaMA2 for illustration, we explore expanding the model’s context window to 128k and 256k tokens through a search procedure. This corresponds to extension ratios of 32$\times$ and 64$\times$, respectively, during this phase.

$\Diamond$ \textit{Fine-tuning to 256k}. Next, the pre-trained LLM undergoes fine-tuning to extend its context window to 256k. To accomplish this, we initially fine-tune LLaMA2 for 400 iterations applying RoPE rescaled factors for 128k. Upon completion, we substitute these factors with the ones corresponding to 256k at the resulting checkpoint, followed by a further 600 steps of fine-tuning. This strategy demonstrates greater efficiency compared to immediately fine-tuning for 256k.

$\Diamond$ \textit{Scaling the fine-tuned extended LLM to 2048k with LongRoPE search}. As a last step, an additional search is applied to the previously fine-tuned LLM at 256k sequence length. Through this approach, we achieve a substantial context window length of 2048k without the need for extra fine-tuning. This process yields a final extension factor of $512\times$.

\textbf{Recovery in shorter context windows}. Upon increasing the context window to an extensive 2048k length, we observe diminished performance within the initial context window range. This degradation aligns with prior observations on positional interpolation~\cite{pi}, which causes position embeddings corresponding to the original context range to be confined to a relatively compact area in the higher-dimensional embedding space. Such compression hampers the language model’s efficacy. In particular, using a 512$\times$ extension factor leads to significant congestion of position embeddings within the foundational 4k context window.

To address this issue, an additional evolutionary search is conducted on the extended LLM, specifically tuning the RoPE rescale factors for shorter context lengths such as 4k and 8k. The upper bound of the searched $\lambda$ is lowered, since shorter sequences necessitate minimal positional interpolation. At inference time, the LLM adaptively modifies the relevant RoPE rescale factors as needed.

\section{LongRoPE}

\label{sec:method}
Building upon these insights, we propose LongRoPE, a method that initially deploys an optimized search strategy to leverage the two identified non-uniformities, and subsequently utilizes this approach to increase the LLM context window to surpass 2 million tokens.

\subsection{Problem Formulation}

These two sources of non-uniformity greatly expand the solution space and complicate the optimization process. To manage this, we recast the multidimensional non-uniform position interpolation optimization as a search problem.

For a LLM targeting a  context window size of $L'$ and lengthy input documents $\mathbf{X}$, where each $\mathbf{x}\in\mathbf{X}$ surpasses $L'$ in token length, we denote the original rotary angle of the $i^{th}$ dimension in RoPE embedding at token position $n$ as  $\frac{n}{\beta^i}$. The optimization problem is then formulated as follows:

\begin{equation}
\begin{gathered}
		\label{eq:problem}

\underset {\mathbf{x}\in\mathbf{X}; \, |\mathbf{x}|\ge L'}{ \operatorname {arg\,min} } \,  \mathcal{L}\, \left( \text{LLM} (\text{RoPE}, \mathbf{X})\right),	\text{with} 
\\
\scalebox{0.93}{
$\underset {\begin{subarray}{l} i=0,\cdot\cdot,\frac{d}{2}-1;\\ n\in[ 0,|\mathbf{x}|);\end{subarray}}{\text{RoPE}_(n)}= {\left[\cdots,\cos\left(\mathbb{I}(\hat{\lambda}_{i}, \hat{n})\frac{n}{\beta^i}\right),  \sin\left(\mathbb{I}(\hat{\lambda}_{i}, \hat{n})\frac{n}{\beta^i}\right),\cdots\right] }$}\\
	and $\mathbb{I}(\hat{\lambda}_{i},\hat{n})=\begin{cases}
	1&\text{ if } n<\hat{n}\\
{\frac{1}{\lambda}_{i}}&\text{ if } n\geq\hat{n}
\end{cases}$
	\end{gathered}
\end{equation}

In this approach, a set of scaling factors $\mathbb{I}(\hat{\lambda}_{i},\hat{n})$ are introduced to address both types of non-uniformity. Here, $\hat{\lambda_i}$ represents the variance in RoPE dimension scaling, and $\hat{n}$ indicates the threshold for token position irregularity. The function $\mathbb{I}(\hat{\lambda}_{i},\hat{n})$ is employed to adjust the rotation angle associated with the $i^{th}$ RoPE dimension; $\hat\lambda_{i}$ acts as the scaling parameter, while $\hat{n}$ serves as the token position cutoff. For token positions preceding $\hat{n}$, the scaling factor $\hat\lambda_{i}$ is not applied, thus the standard RoPE rotary angle $\frac{n}{\beta^i}$ remains unchanged. For positions with $n\ge\hat{n}$, the rescaling is enforced using the introduced factors.

Assuming a target context window length of $L'$, the task is to determine the optimal rescaling coefficients ($\mathbb{I}(\hat{\lambda}_{0},\hat{n})$, $\mathbb{I}(\hat{\lambda}_{1},\hat{n})$, ... $\mathbb{I}(\hat{\lambda}_{i},\hat{n})$, ...) across the $1^{st}$ through $d^{th}$ dimensions of RoPE. The aim is that, with these adjusted $\text{RoPE}$ factors, the target $\text{LLM}$ minimizes the next-token prediction loss $\mathcal{L}$—that is, perplexity—when processing input samples $\mathbf{X}$ of length $L'$.

\subsection{Searching the Non-uniform Position Interpolation}

In order to address the challenge outlined in Eq.~\ref{eq:problem}, we present a straightforward but powerful technique that identifies the best RoPE rescale factors. This approach is designed to maximally leverage the complex and varied structure present across the different dimensions of the position embedding.

\textbf{Search space}. We construct an expansive search space to accommodate the two forms of non-uniformity. The scope of this search space is detailed in Table~\ref{tbl:searchspace}. More precisely, we permit the tuning of an individual rescale factor for each dimension within the RoPE embedding. To streamline the search space formulation, our approach involves optimizing over $\lambda_i$ and $\hat{n}$ rather than directly optimizing $\mathbb{I}(\hat{\lambda}_{i},\hat{n})$, with $\hat{\lambda}_{i}$ defined as $1/\lambda_i$. According to Table~\ref{tbl:searchspace}, $\lambda_i$ can be selected from values ranging from 1.0 (representing simple extrapolation) up to $s\times1.25$ (enabling broader interpolation than PI), using increments of 0.01; here, $s$ corresponds to the ratio of the desired context window expansion.

The parameter $\hat{n}$ determines how many of the first token positions preserve their original RoPE embeddings, foregoing position interpolation. In practice, we permit $\hat{n}$ to be chosen from the set \{0, 1, 2, 4, 8, 12, 16, 20, 24, 28, 32, 64, 128, 256\}. Setting $\hat{n}=0$ means every token position applies the optimized rescale factors.

\textbf{Evolution-based search}. The search space outlined in Table~\ref{tbl:searchspace} encompasses a vast array of positional interpolation methods, which presents substantial difficulties for systematic exploration. For instance, an extension factor of $s=4\times$ results in a combinatorial space of $400^{128/2}\times14=4\times10^{167}$ potential configurations. As the extension ratio increases, the size of the search space grows rapidly in an exponential manner. To manage this complexity, we employ evolution search~\cite{guo2020single} and further implement two optimization strategies that significantly enhance the efficiency of the search process. The overall search framework is described in Algorithm~\ref{alg:search}.

\textit{Optimized initial population generation}. Rather than fully random initialization of the population of $P$ rescale factors, we explicitly include the three RoPE rescale factors associated with PI, NTK, and YaRN as part of the initial set of individuals. The other $P$-3 members of the population are generated by applying random mutations to these three rescale factors, each with a mutation probability $p$.

\textit{Monotonically non-decreasing constraint}. Following the creation of the initial population, we evaluate each individual by measuring the LLM perplexity. In this step, we implement the associated RoPE rescale factors within the target LLM and determine the perplexity for the input $\mathbf{X}$. The best-performing $k$ individuals are then chosen as parents for the next evolutionary phase. Nevertheless, given the high dimensionality of the search space, simple mutation and crossover operations may frequently generate suboptimal candidates, resulting in excessive and redundant perplexity evaluations. This inefficiency becomes more pronounced for larger $L'$, as each perplexity computation is computationally expensive.

To mitigate this issue, we enforce a monotonicity condition where the rescaled RoPE factors sampled must satisfy $\lambda_i\le \lambda_{i+1}$. Only those RoPE configurations meeting this requirement are considered during LLM perplexity assessment, which greatly reduces computational overhead during the search process. More precisely, the rescaling factors $\lambda_i$ are constrained to increase as the RoPE dimension ($i$=0,...,63) grows. This monotonicity across dimensions is informed by NTK theory~\cite{jacot2018neural,tancik2020fourier,ntk}, which posits that lower-indexed dimensions (corresponding to higher frequencies) benefit from minimal interpolation (a smaller $\lambda_i$), whereas dimensions with higher indices (reflecting lower frequencies) are better suited to increased interpolation (a larger $\lambda_i$).

\begin{algorithm}[t]
	\small
	\caption{The search algorithm}
	\textbf{Input:} target LLM, input samples $\mathbf{X}$,   population size $P$, mutation size $N_1$, crossover size $N_2$, max iterations $\mathcal{T}$, mutate probability $p$\\

	\begin{algorithmic}[1]
		\label{alg:search}
		\STATE Top-k$\leftarrow \phi$;
		\STATE $\text{P}_0 \leftarrow$ \textit{Initialize\_population\_with\_optimization} ($P$, $\mathbf{X}$, $p$);
		\FOR{$i=1$ to $\mathcal{T}$}
		\STATE \textit{Compute\_perplexity}(LLM, $\text{P}_{i-1}$, $\mathbf{X}$);
		\STATE Top-k $\leftarrow$ \textit{Update\_Topk}(Top-k, $\text{P}_{i-1}$);
		\STATE $\text{P}_{mutation} \leftarrow$ \textit{Mutation\_with\_mono\_constraint}(Top-k, $N_1$, $p$);
		\STATE $\text{P}_{crossover} \leftarrow$ \textit{Crossover\_with\_mono\_constraint}(Top-k, $N_2$);
		\STATE $\text{P}_i \leftarrow \text{P}_{mutation} \cup \text{P}_{crossover} \cup$ Top-k;
		\ENDFOR
		\STATE Return the member from Top-k exhibiting minimal perplexity;
	\end{algorithmic}
\end{algorithm}

\textbf{8$\times$ extension achieved without fine-tuning}. By utilizing evolutionary search, we successfully determine non-uniform RoPE rescale factors that retain crucial dimensions and positional information, thereby reducing the information degradation from interpolation. As illustrated in Fig.\ref{fig:non-ft}, this approach extends the context window of LLaMA2 from 4k up to 32k without requiring any fine-tuning. Conversely, prior approaches—including PI, and non-uniform NTK and YaRN—experience substantial increases in perplexity after only a 2$\times$ context extension.

\subsection{Extending LLM Context Window to 2048K}
\label{sec:extendto2048k}

\textbf{Progressive extension to 2048k}. In this section, we present our approach for increasing the context window of existing pre-trained LLMs from the standard 4k up to more than 2048k. Our proposed non-uniform positional interpolation enables an extension by a factor of 8 without the need for additional fine-tuning. When even greater context lengths are needed (e.g., scaling by 512$\times$), however, fine-tuning becomes essential. A possible strategy involves optimizing the RoPE rescaling factors for the intended 2048k window and then conducting fine-tuning. This approach, nonetheless, is often hindered by the enormous computational demands involved. Additionally, as indicated by our findings, achieving effective fine-tuning at such large extension scales presents significant difficulties (refer to Appendix).

Luckily, LongRoPE demonstrates strong performance when applied to both base and extended LLMs that have undergone fine-tuning. As such, we propose an efficient and incremental approach, enabling the model to reach the desired 2048k context window using only 1,000 fine-tuning steps, with training conducted on sequences of up to 256k tokens.

$\Diamond$ \textit{Scaling pre-trained LLMs to 256k using LongRoPE search}. Using LLaMA2 as a reference, we explore optimal configurations for context window lengths of 128k and 256k. This results in expansion factors of 32$\times$ and 64$\times$, correspondingly.

$\Diamond$ \textit{Fine-tuning to 256k}. Subsequently, the pre-trained LLM undergoes fine-tuning in order to reach a context window size of 256k. More precisely, we begin by fine-tuning LLaMA2 for 400 steps employing RoPE rescaling factors adjusted for 128k. After completing this stage, we update the checkpoint with RoPE rescaling factors configured for 256k and continue fine-tuning for a further 600 steps. This staged approach demonstrates greater efficiency compared to directly fine-tuning for 256k context length.

$\Diamond$ \textit{Expanding fine-tuned extended LLM up to 2048k using LongRoPE search}. In the final step, we apply an additional search process to the LLM that has already been fine-tuned for 256k sequence length. Through this procedure, the model achieves a remarkably large context window of 2048k, all without requiring any additional fine-tuning. This process yields a total extension factor of $512\times$.

\textbf{Degradation in performance within original context window}. Following the expansion to an exceptionally large 2048k context window, there is an observed decrease in model effectiveness over the standard window length. This limitation is well-documented in the positional interpolation literature~\cite{pi}, where the compression of positional embeddings into a more restricted range within the initial context window impairs the model's ability to encode positional information accurately. When the extension factor reaches 512$\times$, the positional representations corresponding to the original 4k tokens are highly compressed, resulting in increased overlap and reduced distinction between positions.

To address this issue, we conduct an additional evolution search on the expanded LLM, specifically tuning RoPE rescale parameters for brief context spans such as 4k and 8k. The upper limit of the searched $\lambda$ is lowered since shorter contexts demand minimal positional interpolation. In the inference stage, the LLM autonomously modifies its RoPE rescale factors as appropriate.

 \section{Experiments}

\label{sec:exp}
\subsection{Setup}
\textbf{Evaluation Tasks and models}. LongRoPE is integrated with LLaMA2-7B and Mistral-7B, and their capabilities are assessed in three domains: (1) analyzing perplexity scores for long-context LLMs over extended text passages; (2) conducting the Passkey retrieval task, designed to gauge the model's capacity to extract a straightforward passkey amid extensive distractor content; and (3) employing conventional LLM benchmark tasks constrained to a 4096-token context window.

\textbf{Fine-tuning}. For LLaMA2, the model is trained using a linear learning rate decay starting at 2e-5, with a global batch size set to 32. Fine-tuning is conducted for 400 steps using the Redpajama~\cite{redpajama} dataset, which is partitioned into 128k-token segments and encapsulated with BOS and EOS tokens. Upon completion of this stage, a further 600 training steps are performed to expand the context window to 256k tokens, beginning from the obtained checkpoint. The 128k token configuration is trained using 8 A100 GPUs in a distributed manner via the system described in~\cite{cube}, whereas scaling to 256k tokens necessitates 16 A100 GPUs. For Mistral, training employs a fixed learning rate of 1e-6 and a global batch size of 64. Both the 128k and 256k token models adhere to the methodology outlined in YaRN~\cite{yarn}, comprising 400 steps on Together Computer's Long-Data Collections~\cite{mistraldata} with sequences of 16k tokens. The training for Mistral uses a total of 4 A100 GPUs.

\textbf{Search}. For context windows up to 256k tokens, the configuration utilized is: $P$=64, $N_1$=$N_2$=16, $p$=0.3, $\mathcal{T}$=40, and in each round, the top-32 candidates are chosen for mutation and crossover operations. Perplexity evaluation is performed with 5 randomly selected samples from the PG19 validation set, ensuring each meets the minimum context length being tested. When evaluating window sizes greater than 512k, the population size as well as the numbers for mutation and crossover are reduced by half. In this case, perplexity is assessed using 3 randomly drawn samples from the Pile-Books3~\cite{pile} validation set.

\textbf{Baselines}. In order to achieve a context window of 2048k, we fine-tuned models starting with context windows of 128k and 256k. As a result, we obtained LongRoPE-2048k (ft=128k) for LLaMA2 and LongRoPE-2048k (ft=256k) for Mistral. These four variants are benchmarked against leading baselines for extending context windows, specifically focusing on open-source LLMs that have undergone fine-tuning after applying positional interpolation through PI, NTK, and YaRN. The set of baseline models includes Together-32k~\cite{together}, Code LLaMA~\cite{codellama}, LongLoRA-full-FT-100k~\cite{longlora}, as well as YaRN-LLaMA and YaRN-Mistral~\cite{yarn}.

\subsection{Main Results}

\label{sec:mainresult}
\textbf{Long sequence language modeling up to 256k}. Our analysis starts with a comparison against state-of-the-art large language models configured for extended contexts up to 256k tokens. To showcase the generalizability of our approach, we conduct experiments on both the Proof-pile~\cite{pg19} and PG19~\cite{pile} datasets using their respective test splits. Perplexity measurements are taken over various context lengths utilizing a sliding window of 256 tokens. Specifically, for the PG19 dataset, the complete test set consisting of 100 documents is employed. In line with the methodology from YaRN~\cite{yarn}, for Proof-pile we randomly sample 10 examples, ensuring that each selected sample has a minimum length of 128k.

Table~\ref{tbl:proof-pile} and Table~\ref{tbl:pg19} present a comparison of the perplexity scores for LLaMA2 and Mistral, each enhanced using various interpolation strategies, evaluated on the Proof-pile and PG19 datasets, respectively. Two primary findings emerge from these results: \textbf{(1)} the extended models consistently exhibit declining perplexity as the evaluation length increases from 4k to 256k, indicating their enhanced capability to utilize extended contexts; \textbf{(2)} Despite operating with a context window that is 16$\times$ longer—a setting that generally compromises short-context performance—the LongRoPE-2048k models still surpass leading baselines when assessed up to 256k context length.

\textbf{Language modeling for sequences over 2000k tokens}. To assess performance on very lengthy texts, we employ the Books3~\cite{pile} dataset. For practical evaluation, we randomly choose 20 books—each with more than 2048k tokens—and apply a sliding window of 256k tokens.

As illustrated in Table~\ref{tbl:books3}, LongRoPE effectively enlarges the context window of LLaMA2-7B and Mistral-7B models up to 2048k, while maintaining perplexity at a level that is equal to or better than baseline performance for shorter spans in the 8k-128k range. Our results reveal distinct variations between LLaMA2 and Mistral when extended to 2048k contexts. Specifically, Mistral yields superior outcomes compared to baselines at brief context lengths but its perplexity surpasses 7 once the window grows beyond 256k. In contrast, LLaMA2 shows consistent behavior, with perplexity generally decreasing as context length increases, although small rises are noted at 1024k and 2048k. Additionally, LongRoPE-2048k achieves stronger results on LLaMA2 when fine-tuned at 256k rather than 128k, which can be attributed to the reduced secondary extension ratio (8$\times$ versus 16$\times$). Conversely, Mistral demonstrates improved performance when fine-tuned using a 128k window. This outcome stems largely from the training methodology: for Mistral’s 128k and 256k fine-tuning, the YaRN configuration is applied, specifying a 16k training context that restricts Mistral’s potential for further window expansion after fine-tuning.

\textbf{Passkey retrieval}. Next, we investigate how context window size impacts generation performance. To do this, we utilize the passkey retrieval synthetic benchmark described in ~\cite{passkey}. In this benchmark, a model must identify a randomly positioned five-digit number—referred to as the passkey—embedded within an extended document. The structure of the input prompt is specified in the appendix. For evaluation, we conduct 10 rounds of the passkey retrieval task, with the passkey inserted at uniformly random positions throughout the chosen context window.

Fig.~\ref{fig:passkey} illustrates a comparison of retrieval accuracy between our models and the baselines. The performance of baseline models quickly deteriorates to zero once the token count exceeds 128k. Conversely, even when tasked with the formidable challenge of extracting a passkey from millions of tokens, LongRoPE-LLaMA2-2048k (ft=256k) consistently achieves high retrieval accuracy ($\ge$90\%) across a wide range from 4k up to 2048k tokens. Similarly, LongRoPE-Mistral-2048k (ft=128k) maintains perfect accuracy (100\%) until 1800k tokens, after which it declines to 60\% at 2048k. This drop is consistent with observations in Table~\ref{tbl:books3}, where perplexity shows a slight increase at 2048k tokens.

\textbf{Evaluation on standard benchmarks within the initial context window}. LongRoPE-2048k models are assessed on their performance within the default context window using datasets from the Hugging Face Open LLM Leaderboard~\cite{huggingfaceboard}, under both zero-shot and few-shot configurations. Specifically, our evaluation employs a 25-shot setup for ARC-Challenge~\cite{allenai:arc}, a 10-shot configuration for HellaSwag~\cite{zellers2019hellaswag}, a 5-shot regime for MMLU~\cite{mmlu}, and a 0-shot scenario for TruthfulQA~\cite{lin2021truthfulqa}.

As evidenced in Table~\ref{tbl:commonsense}, our models attain results similar to those recorded on the initial benchmark, which was tailored for shorter context lengths. Notably, our models exceed the original Mistral’s performance on TruthfulQA by a margin of +0.5\%. While LongRoPE-LLaMA2-2048k, trained at 256k, exhibits a somewhat greater drop in performance, it still sustains effectiveness across the majority of evaluations.

\subsection{Ablation Results}

\textbf{Effectiveness of the second positional interpolation}. Within the framework of our progressive extension strategy, we implement a secondary non-uniform positional interpolation on extended LLMs using our search algorithm. To assess the impact of this approach, we experiment with our fine-tuned LLaMA2-256k model. This model is further scaled to 512k, 1024k, and 2048k employing both PI and YaRN techniques. According to the results reported in Table~\ref{tbl:secondsearch}, our method of non-uniform positional interpolation maintains perplexity at a stable level as the context window expands. In contrast, perplexity rapidly escalates for both PI and YaRN as the extension size increases.

\textbf{Recovery performance at reduced context lengths.} In order to address the performance degradation observed with shorter contexts, we reconfigure the RoPE factors of LongRoPE-2048k using our search procedure. In this adjustment, we restrict the maximum scale factors in the search, thereby limiting interpolation at 4k and 8k context lengths. Table~\ref{tbl:shortperformance} provides a comparative analysis of perplexity for LongRoPE-LLaMA2-2048k on Proof-pile at both 4k and 8k, as well as reports the mean accuracy across standard LLM benchmarks. The findings illustrate a marked enhancement in performance for brief context lengths.

\textbf{Evaluation of the two types of non-uniformities}. In this section, we investigate the influence of each non-uniformity component through ablation studies. Our approach involves two experimental designs: (i) scaling LLaMA2-7B to shorter context windows (16k and 32k) via several strategies—including PI, optimizing only the RoPE dimension, and jointly optimizing both non-uniformity aspects; (ii) scaling our 256k-length fine-tuned LLaMA2 to an extended length of 2048k using the same methodology. Perplexity metrics are computed prior to any additional fine-tuning. Table~\ref{tbl:searchdimension} highlights that adjusting the RoPE dimension’s non-uniformity yields substantial improvements in perplexity relative to PI's linear interpolation method. The introduction of token position non-uniformity further enhances model performance at 16k and 32k tokens, while its effects at 2048k are less pronounced—potentially because token preservation alone, without interpolation, becomes ineffective at such large context sizes. This phenomenon warrants further exploration and is earmarked as an avenue for future investigation.

\section{Related Works}

Beyond position interpolation methods, this section reviews additional relevant research on alternative strategies.

\textbf{Retrieval-based} approaches utilize an external memory component to store extended historical context, with retrieval mechanisms employed to access pertinent documents during inference~\cite{longllama,wang2023augmenting,borgeaud2022improving}. Such approaches often require substantial architectural changes to large language models. By comparison, our method introduces only minimal changes to positional embeddings, resulting in a more streamlined solution. Furthermore, our technique is adaptable to a wider variety of long-context tasks, including lengthy document summarization and few-shot learning scenarios, rather than being limited solely to retrieval-oriented applications.

\textbf{Attention-based context window extensions}. In addition to positional embedding interpolation techniques, another approach to enlarging the effective context window of LLMs involves adjusting the attention mechanism itself while keeping the native model context length unchanged~\cite{han2023lminfinite, streamingllm,parallelwindow}. The central concept is to address the challenge of attention explosion introduced by extended positions through the use of specialized attention masks. These methods can be utilized in conjunction with positional interpolation, as they provide complementary benefits.

\textbf{Fine-tuning based approaches} investigate strategies for efficiently adapting pre-trained LLMs to accommodate longer contexts through alterations to position embeddings. For instance, methods such as Code LLaMA~\cite{codellama}, LLaMA2 Long~\cite{xiong2023effective}, and ScaledRoPE~\cite{liu2023scaling} select an extremely large RoPE base value and perform fine-tuning on the desired context length. In contrast, our approach enables flexible adaptation to a range of target sequence lengths, successfully scaling beyond 2M tokens. Furthermore, as fine-tuning models for extensive context windows (exceeding 128k tokens) incurs high GPU costs, solutions like LongLoRA~\cite{longlora} and PoSE~\cite{zhu2023pose} have been developed to alleviate these computational demands. Our technique operates independently of such efficient fine-tuning methods.

\section{Conclusion}

This paper introduces LongRoPE, a novel approach that significantly increases the context window of LLMs up to 2048k tokens—an unprecedented scale—while retaining their performance in the original, shorter contexts. Our technique leverages two types of non-uniformities inherent in RoPE positional embeddings, identified through an efficient evolutionary search process. This results in dual advantages: it not only provides optimal initialization for subsequent fine-tuning but also allows for an 8$\times$ expansion of the context window without necessitating fine-tuning. We further develop a progressive extension methodology, whereby models fine-tuned at 256k length are incrementally adapted to achieve a 2048k token context window, all without additional fine-tuning steps. Comprehensive experimental results demonstrate LongRoPE's efficacy. We anticipate that LongRoPE-2048k will unlock a range of novel long-context use cases and stimulate further exploration in this area.

\section*{Broader Impacts} The objective of this study is to contribute to the progression of Machine Learning as a discipline. While the research may carry a range of societal implications, we do not find any particular impact that necessitates explicit mention at this time.

	\balance

\end{document}